% ==============================================================================
\section{MILP Formulation}
\label{sec:milp}
% ==============================================================================

We now present a mixed-integer linear programming formulation of the AHPP
aligned with the aircraft-level view of \S\ref{sec:problem}: each aircraft
$r$ is treated as a single atomic block $[s_r,\, s_r + D_r]$ with total
processing time $D_r$ taken from the instance data, and all blocking
interactions are expressed in terms of the aircraft start and finish
times $s_r, f_r$.  Once an optimal $(s_r^\star, x_{rp}^\star)$ is
obtained, the underlying per-job timings of each aircraft can be
reconstructed deterministically by walking the precedence chain of its
maintenance work, as described in \S\ref{sec:job_recovery}.  The model is
implemented in Pyomo~\citep{Bynum2021} and solved with
Gurobi~\citep{Gurobi2023}.

Throughout this section parameters are written with upper-case Latin
letters and decision variables with lower-case Latin letters, with Greek
letters reserved for fixed scalar constants such as the minimum
separation $\varepsilon$ and the big-$M$ relaxation $M$.

\subsection*{Notation}

\noindent\textit{Sets and indices.}
\begin{tabular}{@{}l@{\quad}p{0.78\textwidth}@{}}
  $\calP$
    & maintenance positions \\
  $\calR$
    & aircraft \\
  $\calA \subseteq \calP \times \calP$
    & blocking arcs $(p, p')$: position $p$ blocks position $p'$ \\
\end{tabular}

\smallskip
\noindent\textit{Parameters.}
\begin{tabular}{@{}l@{\quad}p{0.78\textwidth}@{}}
  $D_r$
    & total processing time of aircraft $r$ (sum of its job durations) \\
  $E_r$
    & earliest start time of aircraft $r$ \\
  $L_r$
    & target finish time (due date) of aircraft $r$; soft, completion
      beyond $L_r$ is penalised but not forbidden \\
  $\varepsilon$
    & tow-in / tow-out time, used both as separation between consecutive
      aircraft at the same position and as the safety margin
      $f_r - \varepsilon$ in the blocking conditions; per-instance value
      read from the input \\
  $H^{UB}$
    & technical upper bound on every aircraft start and finish time,
      $H^{UB} = \max_{r \in \calR} E_r + \sum_{r \in \calR} D_r +
      (|\calR| - 1)\varepsilon$ \\
  $M$
    & big-$M$ constant; we take $M = H^{UB} + \varepsilon$ \\
  $W^M, W^D, W^S$
    & weights for makespan, total delay, and movements \\
\end{tabular}

\smallskip
\noindent
The bound $H^{UB}$ is valid because the schedule that processes every
aircraft sequentially on a single position, starting at $\max_r E_r$ and
inserting a separation of $\varepsilon$ between consecutive aircraft, is
feasible and has makespan exactly $H^{UB}$.  This provides a valid
modelling horizon for the MILP without imposing an operational horizon
on the benchmark.  With $s_r, f_r \in [0, H^{UB}]$, taking $M = H^{UB} +
\varepsilon$ is sufficient to relax the sequencing constraints when
their guard binaries are inactive: the residual right-hand side in
\eqref{eq:sep} (resp.\ \eqref{eq:betaInLB}, \eqref{eq:betaOutLB}) is at
most $f_r + \varepsilon - M \le 0$, so the corresponding constraint
becomes vacuous.  The same value also relaxes the blocking indicator
constraints whenever any of their position or comparison guards are not
active.

\smallskip
\noindent\textit{Decision variables.}
Continuous variables $s_r, f_r \in [0,\, H^{UB}]$ are the start and
finish times of aircraft $r$ at its assigned position.  The assignment
and sequencing binaries are $x_{rp}$ (1 if aircraft $r$ is assigned to
position $p$) and $q_{rr'p}$ (1 if aircraft $r$ precedes aircraft $r'$ at
position $p$).  For every pair of distinct aircraft $r \ne r'$ and every
blocking arc $(p, p') \in \calA$ we add six binary variables that detect
blocking: on the entry side, $\alpha^{\mathrm{in}}_{rr'pp'} = 1$ iff
$s_{r'} > s_r$, $\beta^{\mathrm{in}}_{rr'pp'} = 1$ iff
$s_{r'} \le f_r - \varepsilon$, and $u^{\mathrm{in}}_{rr'pp'} = 1$ iff the
entry of $r'$ at $p'$ forces displacing $r$ at $p$; the exit-side
variables $\alpha^{\mathrm{out}}_{rr'pp'}, \beta^{\mathrm{out}}_{rr'pp'},
u^{\mathrm{out}}_{rr'pp'}$ are defined symmetrically by replacing $s_{r'}$
with $f_{r'}$.  Finally, the makespan $m \in \NR$, the per-aircraft delay
$v^D_r \in \NR$, and the total number of movements $n \in \NZ$ collect the
three components of the objective.

\subsection*{Constraints}

To keep the constraints compact, we abbreviate the universal quantifiers
$\forall r \ne r' \in \calR,\, (p, p') \in \calA$ by $(\dagger)$ and
$\forall r \ne r' \in \calR,\, p \in \calP$ (same-position pair) by
$(\ddagger)$.  The full model is:
\begin{align}
  \sum_{p \in \calP} x_{rp} &= 1
    && \forall r \in \calR \label{eq:assign}\\
  f_r &= s_r + D_r
    && \forall r \in \calR \label{eq:dur}\\
  s_r &\ge E_r
    && \forall r \in \calR \label{eq:earlystart}\\
  s_{r'} &\ge f_r + \varepsilon
              - M(1 - q_{rr'p}) \notag\\
         &\quad - M(1 - x_{rp}) - M(1 - x_{r'p})
    && (\ddagger) \label{eq:sep}\\
  q_{rr'p} + q_{r'rp} &\ge x_{rp} + x_{r'p} - 1
    && \forall r < r' \in \calR,\, p \in \calP \label{eq:ordLB}\\
  q_{rr'p} + q_{r'rp} &\le 1
    && \forall r < r' \in \calR,\, p \in \calP \label{eq:ordUB}\\
  s_{r'} &\ge s_r - M\bigl(1 - \alpha^{\mathrm{in}}_{rr'pp'}\bigr)
    && (\dagger) \label{eq:alphaInLB}\\
  s_{r'} &\le s_r + M\,\alpha^{\mathrm{in}}_{rr'pp'}
    && (\dagger) \label{eq:alphaInUB}\\
  s_{r'} &\le f_r - \varepsilon
              + M\bigl(1 - \beta^{\mathrm{in}}_{rr'pp'}\bigr)
    && (\dagger) \label{eq:betaInUB}\\
  s_{r'} &\ge f_r - M\,\beta^{\mathrm{in}}_{rr'pp'} \notag\\
         &\quad - M(1 - x_{rp}) - M(1 - x_{r'p'})
    && (\dagger) \label{eq:betaInLB}\\
  u^{\mathrm{in}}_{rr'pp'} &\le \min\{x_{rp},\, x_{r'p'},\,
                                   \alpha^{\mathrm{in}}_{rr'pp'},\,
                                   \beta^{\mathrm{in}}_{rr'pp'}\}
    && (\dagger) \label{eq:uInUB}\\
  u^{\mathrm{in}}_{rr'pp'} &\ge x_{rp} + x_{r'p'} \notag\\
                           &\quad + \alpha^{\mathrm{in}}_{rr'pp'}
                           + \beta^{\mathrm{in}}_{rr'pp'} - 3
    && (\dagger) \label{eq:uInLB}\\
  f_{r'} &\ge s_r - M\bigl(1 - \alpha^{\mathrm{out}}_{rr'pp'}\bigr)
    && (\dagger) \label{eq:alphaOutLB}\\
  f_{r'} &\le s_r + M\,\alpha^{\mathrm{out}}_{rr'pp'}
    && (\dagger) \label{eq:alphaOutUB}\\
  f_{r'} &\le f_r - \varepsilon
              + M\bigl(1 - \beta^{\mathrm{out}}_{rr'pp'}\bigr)
    && (\dagger) \label{eq:betaOutUB}\\
  f_{r'} &\ge f_r - M\,\beta^{\mathrm{out}}_{rr'pp'} \notag\\
         &\quad - M(1 - x_{rp}) - M(1 - x_{r'p'})
    && (\dagger) \label{eq:betaOutLB}\\
  u^{\mathrm{out}}_{rr'pp'} &\le \min\{x_{rp},\, x_{r'p'},\,
                                    \alpha^{\mathrm{out}}_{rr'pp'},\,
                                    \beta^{\mathrm{out}}_{rr'pp'}\}
    && (\dagger) \label{eq:uOutUB}\\
  u^{\mathrm{out}}_{rr'pp'} &\ge x_{rp} + x_{r'p'} \notag\\
                            &\quad + \alpha^{\mathrm{out}}_{rr'pp'}
                            + \beta^{\mathrm{out}}_{rr'pp'} - 3
    && (\dagger) \label{eq:uOut}\\
  n &= 2 \!\!\sum_{(p,p') \in \calA}
         \sum_{\substack{r, r' \in \calR\\ r \ne r'}}
         \bigl(u^{\mathrm{in}}_{rr'pp'} + u^{\mathrm{out}}_{rr'pp'}\bigr)
    && \label{eq:movements}\\
  v^D_r &\ge f_r - L_r
    && \forall r \in \calR \label{eq:delay}\\
  m &\ge f_r
    && \forall r \in \calR \label{eq:makespan}
\end{align}

Constraints \eqref{eq:assign}--\eqref{eq:earlystart} pin each aircraft to
a single position, link its finish time to the duration parameter $D_r$,
and enforce the earliest-start window.  The same-position block
\eqref{eq:sep}--\eqref{eq:ordUB} uses the ordering binary $q_{rr'p}$ to
select the disjunction direction whenever two aircraft share a position:
\eqref{eq:sep} forces $s_{r'} \ge f_r + \varepsilon$ when $q_{rr'p} = 1$
and both $r, r'$ are at $p$ (the position guards relax the constraint
otherwise), while \eqref{eq:ordLB}--\eqref{eq:ordUB} guarantee that
exactly one direction is active when the position is shared and that
both directions cannot hold simultaneously.  The entry-blocking block
defines two indicators and combines them:
\eqref{eq:alphaInLB}--\eqref{eq:alphaInUB} make
$\alpha^{\mathrm{in}}_{rr'pp'} = 1$ iff $s_{r'} > s_r$, while
\eqref{eq:betaInUB}--\eqref{eq:betaInLB} make
$\beta^{\mathrm{in}}_{rr'pp'} = 1$ iff $s_{r'}$ falls inside the
servicing window of $r$, with the position guards
$-M(1 - x_{rp}) - M(1 - x_{r'p'})$ in \eqref{eq:betaInLB} switching the
constraint off whenever the blocking placement does not actually hold;
\eqref{eq:uInUB}--\eqref{eq:uInLB} then set $u^{\mathrm{in}}_{rr'pp'}$ to
the conjunction $x_{rp} \wedge x_{r'p'} \wedge
\alpha^{\mathrm{in}}_{rr'pp'} \wedge \beta^{\mathrm{in}}_{rr'pp'}$, so an
entry movement is counted iff both aircraft are at the relevant blocking
positions and $r'$ enters during the servicing window of $r$.  The
exit-blocking block \eqref{eq:alphaOutLB}--\eqref{eq:uOut} mirrors the
entry block with $f_{r'}$ in place of $s_{r'}$.  Finally,
\eqref{eq:movements} aggregates all entry and exit indicators into the
movement count $n$ (the factor $2$ accounts for displacement plus
return), \eqref{eq:delay} lower-bounds the positive delay of each
aircraft by $f_r - L_r$, and \eqref{eq:makespan} lower-bounds the
makespan by every aircraft finish time.

\subsection*{Objective}
\begin{equation}
  \min \quad W^M\, m
         \;+\; W^D \sum_{r \in \calR} v^D_r
         \;+\; W^S\, n.
  \label{eq:milpobj}
\end{equation}

Equation~\eqref{eq:milpobj} aggregates the three operational criteria of
the AHPP into a single scalarised cost.  The first term, $W^M\,m$,
penalises the makespan and therefore rewards compact schedules that free
the hangar early; the second term, $W^D \sum_r v^D_r$, penalises the
total positive delay accumulated past each aircraft's target finish
$L_r$ and is what enforces the soft time windows; the third term,
$W^S\,n$, penalises the number of manoeuvring movements induced by
blocking and is what couples the assignment and scheduling decisions
through the indicators $u^{\mathrm{in}}_{rr'pp'}$ and
$u^{\mathrm{out}}_{rr'pp'}$.  The three non-negative weights
$W^M, W^D, W^S$ are exogenous parameters that the planner can tune to
encode the relative operational priority of throughput, on-time
delivery, and ground-crew effort; \S\ref{sec:exp_setup} evaluates the
model under three contrasting weight profiles.

\subsection{Job-level recovery}
\label{sec:job_recovery}

The atomic-block view of \S\ref{sec:problem} is sufficient for the
optimisation model itself.  In deployment, however, the maintenance work
of each aircraft typically consists of an ordered sequence of jobs
$j^r_1, j^r_2, \ldots$ with individual durations summing to $D_r$.  Given
an aircraft-level solution $(s_r^\star, x_{rp}^\star)$, those per-job
start times are recovered deterministically by walking the chain: the
first job of $r$ starts at $s_r^\star$, and each subsequent job starts at
the finish of its predecessor.  Since none of the constraints above
involve individual job times, this recovery is loss-free: every feasible
aircraft-level solution yields a unique per-job schedule with the same
makespan, total delay, and movement count.
